\documentclass[12pt]{article}
\usepackage[utf8]{inputenc}
\usepackage[T1]{fontenc}
\usepackage{hyperref}
\title{Kuanhua Language Resources for K-12 Curriculum Development}
\author{Compiled by Claude Sonnet 5.5 (claude-sonnet-5-5) for Andr\'as Kornai}
\date{}
\begin{document}
\maketitle

\section{Introduction}
Kuanhua (ISO 639-3 xnh; also Kuan) is reported as a poorly classified Austroasiatic language. A figure of about 1,000 speakers in Jinghong County, Xishuangbanna, Yunnan, goes back to 1991. Its classification is uncertain. One proposal places it in Mangic, and other suggestions include Palaungic, Khmuic, or an independent branch. Glottolog has retired the entry because there is too little information. It suggests Kuanhua may simply be Khmu spoken on the Chinese side of the border. Ethnologue's 2009 data is rated only 20 percent certain. The Glottolog family label in the identification ("Bookkeeping") is a placeholder for such retired entries and is not a real family. The search found no language-specific educational resources.

\section{Primary Educational Resources}
No primers, textbooks, or school programmes for Kuanhua were found.

\section{Linguistic and Reference Materials}
\begin{itemize}
\item Wikipedia article on the Kuan language: \url{https://en.wikipedia.org/wiki/Kuan_language}
\item Glottolog entry (retired): \url{https://glottolog.org/resource/languoid/id/kuan1250}
\item Glottolog entry by ISO code: \url{https://glottolog.org/resource/languoid/iso/xnh}
\item Endangered Languages Project entry and bibliography: \url{https://endangeredlanguages.com/lang/5765} and \url{https://www.endangeredlanguages.com/lang/5765/bibliography}
\item Verbix endangered-languages page: \url{https://docs.verbix.com/EndangeredLanguages/Kuanhua}
\item Related and possibly identical language: Khmu (ISO kjg). Khmu materials may be a more productive route.
\end{itemize}

\section{Religious and Community Resources}
A people-group profile of the Kuan, covering Laos, says they are Mon-Khmer and related to Khmu: \url{https://people-groups.asiaharvest.org/Laos/Kuan-OC.pdf}. It is a missionary-oriented source and should be read with caution. No community organisations were found.

\section{Digital and Audio Resources}
No audio, corpora or software for Kuanhua were found.

\section{Curriculum Development Context}
Evidence is too thin to plan a curriculum. The first step is basic documentation: confirm whether Kuanhua is a distinct language or a Khmu variety, and find speaker communities and local partners in Jinghong. Any teaching would follow community decisions and likely involve Khmu resources.
\end{document}